\chapter{Hurting Yourself (Self-Harm)}

\section*{Definition and Overview}
Self-harm, sometimes called self-injury, is the deliberate act of hurting oneself, most commonly by cutting, burning, or hitting the body. It is usually a way of coping with overwhelming emotional distress rather than an attempt to end life, and it affects people of all ages, particularly adolescents and young adults. Self-harm is a sign of significant distress that deserves compassionate help, not judgement. With the right support, including psychological therapy and treatment of underlying conditions, most people recover and develop healthier ways of coping.

\section*{Epidemiology}
Self-harm is common, affecting around 10--15\% of adolescents and a smaller but significant proportion of adults. Rates have increased in young people, especially girls and young women, and are higher in people with depression, anxiety, trauma, and personality difficulties. Only a minority of self-harm episodes reach medical attention, so true rates are higher than reported figures. Self-harm is a risk factor for suicide, but most people who self-harm do not die by suicide, and early help reduces the risk.

\section*{Aetiology and Pathophysiology}
\begin{itemize}
    \item \textbf{Emotional regulation:} self-harm is often used to release or distract from unbearable feelings
    \item \textbf{Communication:} some people use self-harm to express distress they cannot put into words
    \item \textbf{Self-punishment:} low self-worth and guilt can drive self-harm
    \item \textbf{Trauma:} past abuse and adverse childhood experiences are strongly associated
    \item \textbf{Mental health conditions:} depression, anxiety, borderline personality disorder, and eating disorders
    \item \textbf{Pathophysiology:} self-harm temporarily releases endorphins and provides relief, reinforcing the behaviour as a coping mechanism
    \item \textbf{Contagion:} exposure to self-harm in peers or online can influence behaviour
\end{itemize}

\section*{Clinical Features}
\begin{itemize}
    \item Cutting, scratching, burning, or hitting oneself
    \item Unexplained cuts, scars, or burns, often on the arms, thighs, or torso
    \item Wearing concealing clothing even in warm weather
    \item Withdrawal, irritability, and mood swings
    \item Finding sharp objects or hidden injuries
    \item Difficulties at school, work, or in relationships
    \item The behaviour may be hidden for a long time
\end{itemize}

\section*{Differential Diagnosis}
\begin{itemize}
    \item Accidental injuries, which lack the pattern and secrecy of self-harm
    \item Self-harm must be distinguished from suicidal behaviour, though they can coexist
    \item Underlying conditions: depression, anxiety, trauma, and personality disorders need assessment
    \item Substance use can contribute
    \item A thorough, non-judgemental assessment clarifies the picture
\end{itemize}

\section*{When to Seek Medical Care}
Anyone who is hurting themselves should speak to a GP, who can arrange psychological support and treatment of underlying conditions. Urgent help is needed for serious injury or suicidal intent.

\textbf{Call 000 (or local emergency services) for:}
\begin{itemize}
    \item Serious or severe injury, including deep cuts or burns
    \item Intent to end life or take an overdose
    \item Uncontrolled bleeding
    \item Loss of consciousness
    \item Any situation where there is immediate danger
    \item Call Lifeline (13 11 14) or Kids Helpline (1800 55 1800) for support at any time
\end{itemize}

\section*{Diagnosis and Investigations}
\begin{itemize}
    \item \textbf{Compassionate assessment:} the GP or mental health professional explores the behaviour, triggers, and risk
    \item \textbf{Safety assessment:} the level of risk is assessed
    \item \textbf{Mental health assessment:} screening for depression, anxiety, trauma, and other conditions
    \item \textbf{Medical assessment:} wounds and injuries are examined and treated
    \item \textbf{Referral:} to psychology, psychiatry, or specialist youth services as needed
    \item \textbf{Ongoing review:} follow-up is planned
\end{itemize}

\section*{Management}
\begin{itemize}
    \item \textbf{Psychological therapy:} cognitive behavioural therapy, dialectical behaviour therapy, and other approaches build healthier coping
    \item \textbf{Treat underlying conditions:} depression, anxiety, and trauma are treated
    \item \textbf{Crisis plans:} a written plan of steps to take during urges
    \item \textbf{Safer alternatives:} strategies to manage urges without harm
    \item \textbf{Wound care:} proper treatment of injuries to prevent infection and scarring
    \item \textbf{Family involvement:} education and support for families
    \item \textbf{Online safety:} reducing exposure to self-harm content
    \item \textbf{Medication:} for underlying conditions where appropriate
\end{itemize}

\section*{Complications}
\begin{itemize}
    \item Infection, scarring, and permanent injury
    \item Accidental serious injury or death
    \item Worsening of underlying mental health conditions
    \item Shame, secrecy, and social isolation
    \item Relationship strain and school or work disruption
    \item Increased suicide risk
    \item Financial and legal consequences in rare cases
\end{itemize}

\section*{Prognosis and Outcomes}
Most people who self-harm recover, especially with early, consistent support. Therapy reduces the behaviour in the majority, and most adolescents stop self-harming as they develop other coping skills. Treatment of underlying conditions and a strong support network are the strongest predictors of good outcomes. Recovery takes time and may include setbacks, but self-harm is not a life sentence, and most people go on to live full lives.

\section*{Prevention and Self-Care}
\begin{itemize}
    \item Talk to a GP or trusted adult early, before the behaviour escalates
    \item Engage in psychological therapy and complete treatment
    \item Build a crisis plan and use support helplines
    \item Develop alternative coping strategies: exercise, creative outlets, and grounding techniques
    \item Treat underlying depression, anxiety, and trauma
    \item Reduce exposure to self-harm content online
    \item Family and friends should respond with compassion, not blame
    \item Seek urgent help when there is any suicidal intent
\end{itemize}

\section*{Sources and Further Reading}
The following reputable sources provide further information for healthcare students and patients seeking reliable, current advice about self-harm.

\begin{itemize}
    \item Headspace. \textit{Self-harm support for young people.} https://headspace.org.au/
    \item ReachOut. \textit{Self-harm information and support.} https://au.reachout.com/
    \item SANE Australia. \textit{Self-harm information.} https://www.sane.org/
    \item Kids Helpline. \textit{Support for self-harm.} https://kidshelpline.com.au/
    \item Lifeline. \textit{Crisis support.} https://www.lifeline.org.au/
    \item NHS. \textit{Self-harm: overview and support.} https://www.nhs.uk/mental-health/
\end{itemize}
